\documentclass[a4paper,12pt]{article} % Clase del documento
\usepackage[utf8]{inputenc} % Codificación para caracteres especiales
\usepackage[spanish]{babel} % Soporte para el idioma español
\usepackage{amsmath} % Para escribir ecuaciones matemáticas
\usepackage{amsfonts} % Fuentes matemáticas adicionales
\usepackage{amssymb} % Símbolos matemáticos adicionales
\usepackage{graphicx} % Para incluir imágenes
\usepackage{float} % Para controlar la posición de las figuras (con [H])
\usepackage{geometry} % Ajustar márgenes de la página
\usepackage{hyperref} % Para enlaces clicables (opcional, pero útil)
\geometry{top=2.5cm, bottom=2.5cm, left=3cm, right=3cm}
\title{Análisis de la Calidad del Aire en Medellín utilizando Regresión Lineal por Mínimos Cuadrados}
\author{Carlos Gallego, Gustavo Pérez, Sebastian Pedraza}
\date{\today}

\begin{document}

\maketitle

\section{Análisis de Resultados}

\subsection{Interpretación de los Resultados}

Para estudiar la evolución de la concentración de PM2.5 a lo largo del tiempo, se ajustó un modelo de regresión lineal. La ecuación obtenida es la siguiente:

\begin{equation}
    \text{PM2.5} = m \cdot \text{Days} + b
\end{equation}

donde:
\begin{itemize}
    \item \(\text{Days}\) es la cantidad de días transcurridos desde la fecha más antigua en el conjunto de datos.
    \item \(m\) es la pendiente, que representa la tasa de cambio promedio de PM2.5 por día.
    \item \(b\) es el intercepto o valor inicial estimado cuando \(\text{Days} = 0\).
\end{itemize}

En nuestro caso, los valores específicos obtenidos son:
\[
m \approx 0.012, \quad b \approx 65.345
\]

Para evaluar cuantitativamente el ajuste, se calculó el \textit{Mean Squared Error} (MSE) y el \textit{Root Mean Squared Error} (RMSE):

\[
\text{MSE} \approx 231.69,
\quad
\text{RMSE} \approx 15.22.
\]

Estos resultados indican que, en promedio, el modelo lineal predice la concentración de PM2.5 con un error de alrededor de 15.22 unidades respecto a los valores observados.

\subsection{Discusión de la Calidad del Ajuste y Limitaciones}

La Figura~\ref{fig:pm25_regresion_lineal} muestra los datos reales junto con la línea de regresión. Se observa una gran dispersión en las mediciones de PM2.5 a lo largo del tiempo, lo que sugiere que hay múltiples factores (climáticos, humanos, geográficos, etc.) que influyen en la concentración de contaminantes, y no todos están siendo considerados en este modelo.

Algunas limitaciones son:
\begin{itemize}
    \item \textbf{Variabilidad considerable:} El alto nivel de ruido dificulta que la recta de regresión explique la mayor parte de la variación de los datos.
    \item \textbf{Modelo lineal sencillo:} Es posible que la contaminación siga patrones no lineales o estacionales que un modelo lineal no captura.
    \item \textbf{Falta de variables adicionales:} Factores como tráfico, actividades industriales o parámetros meteorológicos no han sido incluidos y podrían ayudar a mejorar la explicación del fenómeno.
\end{itemize}

\subsection{Propuestas de Mejora o Enfoques Alternativos}

Para mejorar el ajuste y obtener predicciones más realistas, sugerimos:
\begin{itemize}
    \item \textbf{Usar modelos polinómicos o de series de tiempo:} Un modelo de segundo o tercer grado, o un método específico para datos estacionales, podría capturar ciclos o tendencias más complejas.
    \item \textbf{Incluir variables adicionales:} Información climática (temperatura, humedad, velocidad del viento) y de tráfico podrían mejorar significativamente la capacidad predictiva.
    \item \textbf{Aplicar técnicas de Machine Learning:} Algoritmos como Random Forest, Gradient Boosting o incluso Redes Neuronales pueden manejar relaciones más complejas en los datos.
\end{itemize}

\subsection{Visualización de los Resultados}

En la Figura 01 se presenta el gráfico de dispersión de los valores reales de PM2.5 y la línea de ajuste lineal. Además, se incluye en la gráfica el valor del MSE y RMSE para facilitar la interpretación visual:

\begin{figure}[H]
    \centering
    \includegraphics[width=0.75\textwidth]{figures/pm25_regresion_lineal.png}
    \caption{Regresión lineal de la concentración de PM2.5 a lo largo de los días.}
    \label{fig:pm25_regresion_lineal}
\end{figure}

Se evidencia que, pese a la línea de tendencia ligeramente creciente, la alta dispersión de los puntos sugiere la influencia de múltiples factores no controlados por el modelo, limitando su capacidad predictiva.

\subsection{Conclusiones}

En síntesis, la regresión lineal permite observar una tendencia promedio, pero el error cuantificado mediante MSE y RMSE revela que existe gran variabilidad no explicada por el modelo. Por ello, se recomienda emplear enfoques de modelado más sofisticados y la inclusión de variables relevantes para capturar mejor la dinámica de la contaminación del aire en el barrio Belén de Medellín.

\end{document}
